This note supplies the geometric arguments behind Section~5 and
Figures~5--8 of the inspected preprint of Piedade--Tranchida~\cite{PT}.
The constructions are theirs. We use the previously reviewed
component formulation of halving~\cite{Parity}, including its precise
residual-connectivity condition, and the separately reviewed affine
group calculation~\cite{Affine}. In particular a presentation or a
Coxeter diagram will not be used to infer residual connectivity.
All assertions about the source refer to arXiv:2405.19050v1,
not to an uninspected later journal edition.

\section{The cubical model and the retained halving theorem}
Let $n\ge3$, $s\ge2$, and $k\in\{1,2,n\}$. In the usual cubical
tessellation of $\mathbb R^n$ put
\[
 D=\{v\in\Z^n:\textstyle\sum v_i\equiv0\pmod2\},\qquad
 L=D\cap\Lambda_k,
\]
\[
 \Lambda_1=s\Z^n,\quad\Lambda_2=sD,\quad
 \Lambda_n=s(2\Z^n+\Z\one).
\]
A cubical face is a product of integral singleton coordinates and
unit intervals. Its dimension is its number of interval coordinates.
The original toroid has objects the $\Lambda_k$-orbits of faces of
dimensions $0,\ldots,n$, with incidence induced by containment.
We retain the classical assertion that this is the regular polytope
of Theorem~5.1 of~\cite{PT}, attributed there to
\cite[Theorem~6D4]{MS}. Its usual marking, acting on column
coordinates, is
\[
 a=\tau_{e_1}z_1,\quad s_1,\ldots,s_{n-1},\quad z_n,
\]
where $s_i$ swaps coordinates $i,i+1$ and $z_i$ changes the sign of
coordinate $i$. This is the standard affine realization, not a new
classification of cubic toroids.

We recall exactly the halving result used here. A selected vertex--edge
truncation is required to be a simple graph $Q$ (B1). For every other
object, incidence with an edge must be equivalent to incidence with
both endpoints (B2). For a flag $F$ of other objects let $Q_F$ be its
common graph shadow. If $Q$ is nonbipartite, the additional condition is
\[
 \tag{PC}
 \text{all singleton shadows }Q_x\ (x\in F)\text{ nonbipartite}
 \ \Longrightarrow\ Q_F\text{ nonbipartite}.
\]
In the component construction one takes a component of the bipartite
double cover of $Q$, and lifts every other object by the components of
the double cover of its shadow. The resulting incidence system is a
geometry; it is residually connected when $Q$ is bipartite, and, when
$Q$ is nonbipartite, exactly when PC holds. Thinness and chamber
transitivity are preserved. These statements, including the residue
isomorphisms used to compute the diagram, are the retained results of
\cite{Parity}. The regular coset identification and full type-preserving
automorphism group are supplied by~\cite{Coset}. Neither result permits
one to omit PC in a nonbipartite step.

Write $H_{ab}$ for the geometry after $a$ operations at the left end
and $b$ at the right end, with $H_{00}$ the original toroid. The notation
will only be used for the sequences proved to exist below. Choices of
the two equivalent leaves at a fork are identified by a type
permutation. A ``second operation at the left'' means the leaf
$(0,2)$ after the first operation $(0,1)$.

\section{Small translations and local cubical incidence}
The following elementary separation facts keep the quotient arguments
explicit. Put $R=\{\pm e_i\pm e_j:i\ne j\}$.

\begin{lemma}\label{NI.separation}
Every nonzero vector of $\Lambda_k$ or $L$ has a coordinate of absolute
value at least $s$. No two vertices of a unit cube are therefore
identified. Also:
\begin{enumerate}
\item $2e_i\in\Lambda_k$ or $2e_i\in L$ occurs exactly when
$(k,s)=(1,2)$.
\item If $r,r'\in R$ are distinct and $r-r'\in L$, then
$(k,s)=(1,2)$ or $(k,s)=(2,2)$. Conversely both parameter pairs
allow such a coincidence.
\item Suppose neither of these two parameter pairs occurs. If two
even vertices of one unit cube differ by a root modulo $L$, they
differ by that root literally. Their unique square is a face of
the cube. If two even vertices at root distance have a common odd
unit-step neighbor modulo $L$, that neighbor is one of the other
two corners of this square, after a common lift.
\end{enumerate}
The corresponding statements hold in the dual cubical grid of
cube centers, translated by $\frac12\one$.
\end{lemma}
\begin{proof}
All lattices are contained in $s\Z^n$. For (1), it remains only to
check $s=2$. The vector $e_i$ is neither in $D$ nor in
$2\Z^n+\Z\one$ for $n\ge3$.

For (2), every coordinate of $r-r'$ has absolute value at most two.
Thus $s=2$. If $r-r'$ is even coordinatewise, the supports of $r,r'$
are equal, and its nonzero support has size at most two. In $2D$ a
single coordinate $\pm2$ is excluded; a pair $\pm2,\pm2$ is allowed.
In $2(2\Z^n+\Z\one)$ neither a one-coordinate nor a two-coordinate
vector of this form occurs, since $n\ge3$. Passing to $L$ cannot
create a coincidence. Examples are $e_1+e_2$ and $-e_1+e_2$ for
$(1,2)$, and $e_1+e_2$ and $-e_1-e_2$ for $(2,2)$.

For (3), lift the cube and its two vertices $p,q$. In
$q-p-r\in L$, every coordinate has absolute value at most two.
For $s\ge3$ it is zero. For $s=2$, a nonzero even coordinate
can occur only on the two coordinates supporting $r$, and the
remaining allowed lattice $2(2\Z^n+\Z\one)$ contains no such
nonzero vector. Hence $q-p=r$. The square has exactly these two
varying coordinates, so is unique and contained in the cube.

For the common-neighbor assertion, first lift the neighbor $w$ and
one endpoint $p$ so that $w-p$ is a unit vector. The other incidence
gives $q-p\equiv\alpha+\beta\pmod L$, where $\alpha,\beta$ are
signed unit vectors. This sum is zero, a root, or $\pm2e_i$.
The zero case cannot equal a root modulo $L$. Two roots are equal
by (2). Finally $\pm2e_i-r$ cannot belong to $s\Z^n$ for any
$s\ge2$: at least one coordinate belonging to the other coordinate
of $r$ has value $\pm1$. Thus the sum is the same root, and the
two possible intermediate vertices are exactly the other square
corners. Translation of the entire argument gives the dual statement.
\end{proof}

We will also use the following form of local lifting. Fix a unit cube
$C$, or a cube in the dual grid. Its vertices and subfaces inject into
the quotient: a nonzero period cannot move a vertex or a subface of
$C$ to another vertex or subface of $C$, since the coordinate
differences have absolute value at most one. An incident chain of
face orbits has a common lift by first fixing its largest face and
then lifting the smaller faces inside it. These lifts are unique.
Incidences between subfaces of that largest face are consequently
literal containments. When a vertex and a chain of faces form a flag,
the vertex has its unique lift inside the smallest face. These
observations apply to primal and dual cubes, and will be used together
with Lemma~\ref{NI.separation} when a flag has two distinct vertex types.

\begin{proposition}\label{NI.first}
The original left and right leaves satisfy B1 and B2 exactly when
$(k,s)\ne(1,2)$. After the first left operation the geometry $H_{10}$
has the following literal model modulo $L$:
\[
\begin{array}{c|l}
0&\text{even integral vertices}\\
1&\text{odd integral vertices}\\
i\ge2&\text{original }i\text{-cubes}.
\end{array}
\]
Incidence between types 0 and 1 is unit-step adjacency; all incidences
with cubes and between cubes are containment. This geometry is a
regular hypertope.
\end{proposition}
\begin{proof}
Unit-step edges can become parallel in the quotient only if two
distinct signed unit vectors differ by a period. Because periods are
in $s\Z^n$, this is exactly the case $2e_i\in\Lambda_k$, settled
by Lemma~\ref{NI.separation}(1). To prove B2, lift a cube and two of its
vertices. If their difference is a unit vector modulo $\Lambda_k$,
the error has coordinates of absolute value at most two and can have
a nonzero even coordinate only on that single axis. The only possible
period is again $\pm2e_i$ at $(k,s)=(1,2)$. Otherwise the edge is a
literal edge of the cube. The forward implication is containment.
Self-duality proves the right-end assertion.

The parity of the coordinate sum colors the infinite grid. Its
color-preserving period subgroup is exactly $L$. The component
construction on the quotient by $\Lambda_k$ is consequently the
displayed colored-vertex construction modulo $L$. More explicitly,
if $\Lambda_k\subset D$ the quotient grid is bipartite and no
objects are doubled; if an odd period exists, there are two cosets
of $L$ in $\Lambda_k$, giving the two vertex lifts and the two
compatible lifts of each cube. The three displayed lattices lie in
$D$ exactly when $ks$ is even; an odd period otherwise supplies
an odd closed walk in the quotient grid. Every cube shadow is an ordinary
cube graph and is bipartite; higher flags are chains of cubes and
their shadows are the smallest cubes. Thus PC holds, including
in the globally nonbipartite case. The retained halving theorem
gives thinness, residual connectivity and regularity.
\end{proof}

\section{A second operation at one end}
\begin{proposition}\label{NI.second}
In $H_{10}$ the leaves $(0,2)$ and $(1,2)$ satisfy B1 and B2
exactly when $(k,s)\ne(2,2)$, assuming $(k,s)\ne(1,2)$ already.
Under this condition the second-left graph is nonbipartite, PC
holds, and its halving $H_{20}$ is a regular hypertope. Its objects
modulo $L$ are
\[
\begin{array}{c|l}
0,2&\text{two separate copies of the even integral vertices}\\
1&\text{odd integral vertices}\\
i\ge3&\text{original }i\text{-cubes}.
\end{array}
\]
The two even-vertex types are incident when their underlying vertices
differ by a root in $R$; even and odd types are incident at unit
distance; incidences with cubes and between cubes are containment.
\end{proposition}
\begin{proof}
Each square has exactly two even corners. Its resulting edge has
a root as difference vector. Conversely a root specifies a unique
square. Lemma~\ref{NI.separation}(2) therefore proves simplicity except
at the stated small period. For B2 at a higher cube use part (3) of
that lemma. At an odd vertex $w$, the even neighbors are the $2n$
points $w\pm e_i$. They are distinct by part (1). Two such neighbors
give a square through $w$ exactly when they are on different axes.
Two opposite neighbors on the same axis cannot differ by a root
modulo $L$, by the last argument in part (3). This proves B2.

At $(k,s)=(2,2)$ the squares on coordinates 1,2 with respective
corner sets $\{0,e_1,e_2,e_1+e_2\}$ and
$\{0,-e_1,-e_2,-e_1-e_2\}$ are distinct, but have the same
two even corners modulo $2D$. Hence B1 fails. B2 fails too:
the odd vertex $e_1$ is adjacent to both these even corners but
is not a corner of the second square modulo $2D$.

The root graph is connected because $R$ spans $D$, and contains
the triangle $0,e_1+e_2,e_1+e_3$. Its three vertices remain
distinct. The graph shadow of an $m$-cube, $m\ge3$, is its
demicube: one parity class of its corners, adjacent at Hamming
distance two. It is connected and contains a triangle. Indeed,
any even difference can be removed two coordinates at a time,
and the three corners with difference supports
$\varnothing,\{1,2\},\{1,3\}$ form a triangle (with a translated
choice for the other parity). The shadow at an odd vertex is the
cross-polytope graph on $\{w\pm e_i\}$, with edges between
nonopposite pairs. It is connected and contains a triangle for
$n\ge3$.

A higher flag is either a chain of cubes, whose common shadow is
the demicube of its smallest member, or such a chain together
with an odd vertex $w$. In the latter case its even neighbors
inside the smallest $m$-cube are the $m$ corners obtained from
$w$ by a single allowed step; they form $K_m$, with $m\ge3$.
If there is no cube the shadow is the full cross-polytope graph.
Local lifting justifies these descriptions in the quotient. All
of them are nonbipartite, proving PC. Thus every higher object has
just one lift in the new component construction, while the even
vertices acquire two copies. This gives precisely the asserted
model and proves its regular-hypertope conclusion.
\end{proof}

There is a useful consequence of this proof. Any flag of just the
three lower types in $H_{20}$ has a common lift whose vertices span
an original face of dimension at most two. A single vertex spans a
point; an even--odd pair spans an edge; a pair of even types spans
their unique square. If the third, odd, type occurs, Lemma~\ref{NI.separation}
puts it at one of the other corners of that square. The span is
unique. If a cube is also prescribed, the span lies inside its
unique lift. For $H_{10}$ the corresponding lower flags span at
most an edge. We call this the \emph{small-face property} below.

\section{Operations at the opposite end}
Cube centers form the dual grid $\frac12\one+\Z^n$. Color a
center $c=\frac12\one+v$ by the parity of $\sum v_i$.
This is well-defined modulo $L$. For a primal $d$-face $F$, the
centers of the $n$-cubes containing $F$ form the vertices of a
dual cube of dimension $n-d$. A facet between two such cubes
is a dual edge. These are literal local statements, also after
quotienting by $L$, by the local-lifting observation.

\begin{proposition}\label{NI.oppositefirst}
One may perform the first right operation after $H_{10}$ for all
$n\ge3$, and after $H_{20}$ for $n\ge4$. The selected graph is
bipartite, B1 and B2 hold, and the resulting geometries $H_{11}$
and $H_{21}$ are regular hypertopes.
Their two upper types $n,n-1$ are respectively the two colors of
original $n$-cube centers. These types are incident precisely when
the cubes share a facet. Lower objects retain their earlier model,
and are incident with an upper cube when they are contained in it.
The cubes previously of dimension $n-1$ are removed.
\end{proposition}
\begin{proof}
In both cases the original $n$-cubes and $(n-1)$-cubes are still
present: for $H_{20}$ this uses $n-1\ge3$. Their graph is the
dual unit-step grid modulo $L$. It is simple by
Lemma~\ref{NI.separation}(1), and bipartite by the displayed coloring.
B2 for a lower original cube is the primal/dual version of the
edge argument in Proposition~\ref{NI.first}. For a lower vertex,
including either copy of an even vertex, it is the same argument
for the dual $n$-cube around that vertex. No other object types
occur. Every higher shadow is a connected bipartite cube graph:
for a flag without an original cube, use its small-face span;
when there is a chain of original cubes, use the largest one.
Their codimension is at least two. The component construction
therefore splits the center graph into its two colors, does not
duplicate any other object, and has precisely the model asserted.
Residual connectivity follows from the bipartite case of the
retained theorem, without an additional PC assumption.
\end{proof}

\begin{proposition}\label{NI.oppositesecond}
A second right operation is possible in $H_{11}$ when $n\ge4$
and in $H_{21}$ when $n\ge5$, provided $(k,s)\ne(2,2)$.
Its graph is nonbipartite, B1 and B2 hold, and PC holds. It
therefore gives $H_{12}$ (dual to $H_{21}$) and $H_{22}$,
respectively, as regular hypertopes.
The two new upper types $n,n-2$ are separate copies of the even
centers, incident at root distance; type $n-1$ consists of odd
centers. All remaining objects keep their earlier meanings.
\end{proposition}
\begin{proof}
The objects of type $n-2$ are still original $(n-2)$-faces under
the rank assumptions. They have four incident $n$-cubes, two
with even centers. Thus the selected graph is the root graph
on the even dual vertices. The proofs of simplicity and B2
are exactly Lemma~\ref{NI.separation} in the dual grid, applied
separately to each actual higher object: a remaining original
face has a dual cube as its shadow, a lower vertex has a dual
$n$-cube, and an odd center has its cross-polytope neighbor graph.
Different copies of a lower vertex have the same singleton shadow,
which causes no problem for B2.

For completeness consider an arbitrary higher flag, not just its
singletons. First suppose it contains no odd center. If it contains
original faces, their largest member $F$ determines the common
dual cube, since all the lower vertices are in $F$. Its dimension
is $n-\dim F\ge3$, because type $n-2$ is being halved and
the remaining face dimensions are at most $n-3$.
If the flag contains only lower vertex types, their small-face
span has dimension $d\le1$ in $H_{11}$ and $d\le2$ in $H_{21}$.
The common even-center shadow is the demicube of a dual cube
of dimension $n-d\ge3$. This also covers a singleton. An empty
flag has the whole root graph, which is connected and nonbipartite.

Now suppose an odd center $c$ is in the flag. Lift its primal
$n$-cube. Every prescribed lower vertex and every prescribed
primal face lifts inside this cube. The small-face property and
local lifting make these lifts mutually compatible. Let $F$ be
their largest face, or their small-face span if there is no
prescribed face. A neighboring even-center cube contains $F$
exactly when its common facet with the cube at $c$ contains $F$.
There are $n-\dim F$ such facets, so the shadow is the clique
$K_{n-\dim F}$. Its size is at least three by the same bounds.
If there are no other objects, the shadow is the full
cross-polytope graph. All these graphs are connected and
nonbipartite. This proves PC for every flag. The component
construction now gives exactly the stated duplicate even-center
model. The retained theorem proves the conclusion.
\end{proof}

These arguments also identify the order of the operations. Where
both are defined, doing the first right operation before the second
left operation yields the same $H_{21}$ incidence system. One can
either apply Proposition~\ref{NI.oppositesecond} by duality to $H_{11}$,
or compare the displayed models: both have two even-vertex copies,
odd vertices, the remaining primal cubes, and the two center colors,
with exactly the incidences already specified. The same comparison
identifies the two paths to $H_{22}$. Thus an arrow in a diagram of
constructions does not rely merely on equality of the group orders.

\section{Exact diagrams and exhaustion of leaf operations}
The regular-marking theorem in~\cite{Coset} replaces a selected leaf
generator $r$ by $rqr$, where $q$ is its neighbor. Therefore the
markings of the models just proved have the following uniform form:
\[
 (\ell_a,s_1,\ldots,s_{n-1},u_b),
\]
where
\[
\begin{array}{c|c|c}
j&\ell_j&u_j\\\hline
0&\tau_{e_1}z_1&z_n\\
1&\tau_{e_1+e_2}z_1z_2(1\,2)&z_{n-1}z_n(n-1\,n)\\
2&\tau_{e_1+e_3}z_1z_3(1\,3)&z_{n-2}z_n(n-2\,n).
\end{array}
\]
Here $(i\,j)$ denotes a coordinate swap. The second row is the
first conjugation at each end; conjugating the next inward swap
gives the third. This also checks all overlap cases in small rank.
Since the geometries are thin and residually connected, these
product orders are the actual rank-two residue gonalities.

We give indexed edge tables to make the small-rank diagrams
unambiguous. Unlisted pairs have order two. Let
\[
 P=\{\{i,i+1\}:1\le i\le n-2\}.
\]
Its edges all have label three. In the next table, an edge listed
in the last column has label four; all other listed edges have
label three. Repetitions in a union do not add another edge.
\[
\begin{array}{c|l|l}
\text{state}&\text{edges of label three}&\text{edges of label four}\\\hline
00&P&01,\ (n-1)n\\
10&P\cup\{02\}&(n-1)n\\
11,\ n=3&P\cup\{02,03,13\}&\varnothing\\
11,\ n\ge4&P\cup\{02,(n-2)n\}&\varnothing\\
20,\ n=3&P\cup\{01,02\}&03,23\\
20,\ n\ge4&P\cup\{01,02,03\}&(n-1)n\\
21,\ n=4&P\cup\{01,02,03,04,24\}&\varnothing\\
21,\ n\ge5&P\cup\{01,02,03,(n-2)n\}&\varnothing\\
22,\ n=5&P\cup\{01,02,03,05,25,35,45\}&\varnothing\\
22,\ n\ge6&P\cup\{01,02,03,(n-3)n,(n-2)n,(n-1)n\}
 &\varnothing
\end{array}
\]
For example, $02$ means the unordered pair $\{0,2\}$.
The states $21$ and $22$ only occur at the ranks in the table.

Here is a direct verification of the product orders, independent of
any drawn diagram. Interior swaps on disjoint supports commute;
consecutive ones have product order three. A signed swap on
coordinates $i,j$ has product order three with a swap sharing
exactly one coordinate, and commutes with a disjoint swap. The
two opposite signed swaps on the same pair commute. A sign change
on one coordinate of a signed swap has product order four.
The same statements apply to the affine left reflections after
translating their common fixed intersection to the origin. In the
few overlapping left--right pairs one can check the displayed
maps directly: for $n=3$, $\ell_1u_1$ has order three and
$\ell_2z_3$ has order four;
for $n=4$, $\ell_2$ and $u_1$ have order three; for $n=5$,
$\ell_2$ and $u_2$ have order three. The linear parts already
have these orders. Thus the orders cannot decrease after a
translation quotient. Pairs declared commuting are commuting
as affine maps, so the table holds for every allowed period.

\begin{theorem}\label{NI.complete}
Exclude $(k,s)=(1,2)$. Up to duality and exchange of the two
equivalent fork leaves, the complete sequence of admissible leaf
halvings has states
\[
 00\longrightarrow10\longrightarrow11,
\]
and, if also $(k,s)\ne(2,2)$,
\[
 10\longrightarrow20,\qquad
 20\longrightarrow21\longleftarrow11\quad(n\ge4),\qquad
 21\longrightarrow22\quad(n\ge5).
\]
The initial arrow uses the bipartition construction if $ks$ is even
and the partition construction if $ks$ is odd. The arrows at the
first opposite end are bipartition operations. Each second operation
at an end is a nonbipartite partition operation.
For $n=3$, states $20$ and $11$ are terminal; for $n=4$, state
$21$ is terminal; for $n\ge5$, state $22$ is terminal.
Every permitted state is a regular hypertope, with the exact diagram
in the table.
\end{theorem}
\begin{proof}
Existence, nondegeneracy, parity and PC for all arrows have been
proved in Propositions~\ref{NI.first}--\ref{NI.oppositesecond} and their
dual forms. The original left endpoint reflection $a$ exchanges
vertex parity and the two left fork markings. At the other end
$z_n$ preserves vertex parity and exchanges center parity; it
conjugates the two right fork markings. Whenever a second operation
at the opposite end is already present, the stated rank bounds
make these endpoint reflections commute with its altered marker.
Thus these equivalences still apply at that stage.
The explicit models identify both orders of independent operations.

It remains to justify exhaustion. The table shows that $00$ has
only its two endpoints as leaves. State $10$ has its two left
fork leaves and its right endpoint. At $(k,s)=(2,2)$ the former
fail B1 by Proposition~\ref{NI.second}; the latter still works.
State $20$ has no leaf when $n=3$, and otherwise only the right
endpoint. State $11$ has no leaf when $n=3$; otherwise its four
fork leaves lead, up to the symmetries already described, to $21$.
State $21$ has no leaf when $n=4$, and otherwise just the two
right fork leaves. Finally every vertex in state $22$ has degree
at least two. These are exactly the states and arrows stated.
At the excluded second-step period, the same dual root coincidence
prevents the fork operations in $11$ as well. There is therefore
no additional leaf operation omitted from the classification.
\end{proof}

This is the precise interpretation of the families in Figures~5--8:
the second-step arrows require the extra small-period exclusion;
the last two graphs are the overlapping-end cases in ranks four
and five, not disjoint copies of the large-rank picture. The table
specifies every bond even where a reproduced figure is difficult
to read. The proof supplies the nondegeneracy and PC checks which
the final paragraph of Section~5.3 leaves implicit. It does not
restore the insufficient exact presentations in Propositions
5.8--5.13: the corrected marked groups and all missing period
relations are treated separately in~\cite{Affine}.

\section{Scope and attribution}
The toroids, the halving operation, the proposed diagram families,
and their attribution to earlier polytope and diagram geometry are
retained from the source. The new argument here is an explicit
coordinate verification of every allowed geometric branch, using
the previously reviewed parity-shadow criterion rather than the
source's overbroad residual-connectivity claim. It is a Section~5
recovery component. It makes no whole-original admission, numerical
compression claim, new conjecture claim or formalization claim.
